\section{DBSCAN：通过核心点扩展密度连接}
\label{sec:unsup-dbscan}
DBSCAN是基于密度的聚类，不预先指定簇数。
给定距离阈值$\varepsilon$与最少邻域点数$m$，本节约定邻域包含点自身：
$N_\varepsilon(x)=\{x_j:d(x,x_j)\leq\varepsilon\}$。

\begin{figure}[htbp]
\centering\begin{tikzpicture}[x=4cm,y=1cm]
\draw[Quiet,->] (-.15,0)--(3,0) node[right,lab] {$x$};
\draw[Teal,dashed] (.04,.85)--(.56,.85);
\draw[Teal] (.04,.73)--(.04,.97);\draw[Teal] (.56,.73)--(.56,.97);
\node[lab,above] at (.3,.95) {$0.3\pm\varepsilon$};
\foreach \x in {.1,.2,.3,2.6,2.7}{\fill[Teal] (\x,0) circle(2.7pt);}
\foreach \x in {0,.55,2.5,2.8}{\draw[Amber,thick,fill=white] (\x,0) circle(2.7pt);}
\draw[Ink,thick] (1.48,-.08)--(1.52,.08) (1.48,.08)--(1.52,-.08);
\foreach \x in {0,.1,.2,.3,.55,1.5,2.5,2.6,2.7,2.8}{\node[lab,below,rotate=55,anchor=north east] at (\x,-.07) {\x};}
\node[lab] at (1.4,-1.4) {实心：核心点\quad 空心：边界点\quad 叉号：噪声};
\end{tikzpicture}

\caption{一维例子，$\varepsilon=0.26,m=4$。点的类型由邻域计数决定。}
\par\smallskip
\readfigure{0.3的邻域为$\{0.1,0.2,0.3,0.55\}$，因此是核心点；
0.55仅有两个邻域点，是靠近核心点的边界点；1.5既非核心，也不邻接任何核心，是噪声。}
\end{figure}

\begin{center}
\begin{tabularx}{\textwidth}{P{28mm}Y}
\toprule
点的类型 & 判定与扩展\\
\midrule
核心点 & $|N_\varepsilon(x)|\geq m$，可继续扩展邻域。\\
边界点 & 不满足核心条件，但落入某核心点的邻域；可以归簇，不能继续作为扩展桥梁。\\
噪声点 & 没有归入任何密度簇的点；“噪声”不自动意味着测量错误。\\
\bottomrule
\end{tabularx}
\end{center}

先连接彼此邻接的核心点，再吸收邻近边界点，可以得到两个簇，1.5留作噪声。
边界点若同时邻接不同簇，归属可能受遍历约定影响；不能通过一串非核心边界点任意合并簇。

尺度和距离仍很关键：$\varepsilon$过小可能几乎全为噪声，过大可能合为一簇。
不同密度的簇未必能用同一阈值识别；高维距离也会变得困难。
“形状灵活”不等于对数据没有结构假设。

\section{层次聚类：链接方式决定合并距离}
\label{sec:unsup-hierarchy}
凝聚式层次聚类（Agglomerative Clustering）从每点一簇开始，
每次合并距离最近的两个簇。树状图（Dendrogram）的高度表示合并距离，
切割高度可产生不同簇数。

\begin{figure}[htbp]
\centering\begin{tikzpicture}[x=1cm,y=.8cm]
\foreach \panelshift/\mergeheight/\linkname in {0/3/单链接,7/6/全链接}{
\begin{scope}[xshift=\panelshift cm]
\draw[Quiet,->] (-.6,0)--(-.6,6.5) node[above,lab] {距离};
\foreach \y in {1,2,3,6}{\node[lab,left] at (-.6,\y) {\y};}
\draw[Teal,thick] (0,0)--(0,1)--(1.3,1)--(1.3,0);
\draw[Teal,thick] (2.6,0)--(2.6,2)--(3.9,2)--(3.9,0);
\draw[Teal,thick] (.65,1)--(.65,\mergeheight)--(3.25,\mergeheight)--(3.25,2);
\foreach \x/\label in {0/0,1.3/1,2.6/4,3.9/6}{\node[lab,below] at (\x,0) {\label};}
\node[lab,below] at (1.95,-.65) {\linkname};
\end{scope}}
\end{tikzpicture}

\caption{同样四点0、1、4、6，链接方式不同，最后合并高度不同。}
\par\smallskip
\readfigure{先在距离1处合并0与1，再在距离2处合并4与6。
随后两簇的单链接距离为3，全链接距离为6；平均链接为4.5。簇内点未移动。}
\end{figure}

单链接取跨簇点对的最小距离，容易形成链；全链接取最大距离；
平均链接取跨簇点对距离的平均。
Ward方法则以组内平方和增加为准则，常用于欧氏设置，不能混同于上述最短/最长距离。

\quick{聚类都需要选表示、尺度、距离与规则。比较结果时说明簇数、噪声处理和用途，不能只按图形好看判断。}

\begingroup\small
\noindent\textbf{参考：}Ester等，DBSCAN（1996）；ISLR第二版第12章；scikit-learn Clustering文档。
\par\endgroup
